% TOP_PROBLEM: {"id": 30007261, "title": "Odd regular Fontaine-Mazur over Q at two", "category_id": 1, "category": {"id": 1, "name": "number_theory", "display_name": "Number Theory", "description": "Properties of integers, prime numbers, Diophantine equations.", "slug": "number-theory", "order_index": 1, "created_at": "2026-07-31T15:26:25.670Z"}, "status": "open"}
% FIRST_POSTED: 2026-10-06
% ATTEMPT_STATUS: solved
% OPENAI_CLAIM: {"schema_version":1,"source_repo":"https://github.com/openai/math","source_commit":"adc7f1241b42e322a6451854ab7e4b4c146bf78a","release_date":"2026-10-06","match":"full","resolution":"proved","independently_reviewed":false,"second_pass":"agreed","families":[{"family":"010","title":"Unrestricted pro-modularity at the prime two","subject":"Number theory","manuscripts":[{"title":"Fontaine–Mazur modularity at the prime 2","pdf_path":"preprints/Fontaine-Mazur-modularity-at-the-prime-2-September-23-2026/paper.pdf","readme_path":"preprints/Fontaine-Mazur-modularity-at-the-prime-2-September-23-2026/README.md","date":"2026-09-23","theorem_ref":"thm:main (sections/introduction.tex:20)","statement_sources":[{"path":"preprints/Fontaine-Mazur-modularity-at-the-prime-2-September-23-2026/build/sections/introduction.tex","line":20,"label":"thm:main"}]}],"lean":null,"reasoning_trace":null}]}
\documentclass[11pt]{article}
\usepackage[margin=25mm]{geometry}
\usepackage{fontspec}
\setmainfont{Latin Modern Roman}
\usepackage{amsmath,amssymb,mathtools}
\usepackage{unicode-math}
\setmathfont{Latin Modern Math}
\usepackage{enumitem}
\usepackage{xurl,hyperref,bookmark}
\hypersetup{colorlinks=true,urlcolor=blue}
\setcounter{secnumdepth}{0}
\setlength{\parindent}{0pt}
\setlength{\parskip}{5pt}
\setlength{\emergencystretch}{3em}
\begin{document}
\section*{Odd regular Fontaine-Mazur over Q at two}

\subsection{OpenAI's claim}
OpenAI presents the following claims relevant to this question, with manuscript numbers following paper order in Sources. In manuscript 1 (family 010; 2026-09-23), the asserted result is that an odd irreducible two-dimensional two-adic representation over the rationals, finitely ramified and de Rham with unequal Hodge--Tate weights, comes from a classical modular form after a Tate twist.

\subsection{Mathematical statement of the claimed result}
\paragraph{Manuscript 1.}
Let $r:G_{\mathbb Q}\to\operatorname{GL}_2(\overline{\mathbb Q}_2)$ be continuous and irreducible. Assume it is odd, unramified outside a finite prime set, and has a de Rham restriction at $2$ whose two Hodge--Tate weights differ. The claimed modularity conclusion is
\[
 r\simeq r_{f,2}(m)
\]
for some classical cuspidal eigenform $f$ and integer Tate twist $m$. No residual irreducibility or prescribed local ordinary condition is part of the displayed theorem.

\subsection{Outline of the argument}
For manuscript 1 (family 010; 2026-09-23), local deformation charts and solvable base change create a family with modular support. Integral completed forms and their block action permit patching, after which an ordinary vector and the prescribed local line give classicality and descent.

\subsection{Scope relative to this problem}
The result supplies the conclusion in this problem with its displayed hypotheses and quantifiers. The broader statements and applications in the manuscript are kept separate from the particular question above. The supporting result is manuscript 1 (family 010; 2026-09-23).

\subsection{Formal verification}
No Lean formalization is listed at this snapshot for manuscript 1 (family 010; 2026-09-23).

This site has not executed the Lean code. Described statement coverage is not an independent proof check.

\subsection{Status caveat}
The manuscripts were produced by an internal OpenAI model. OpenAI cautions that results without formalization can contain mistakes. This site has not independently checked the proofs; the label records the scope of OpenAI's claim.

\subsection{Sources}
\begin{itemize}
\item OpenAI, \emph{Fontaine–Mazur modularity at the prime 2}, 2026-09-23; citation key \texttt{OAI:Fontaine-Mazur-modularity-at-the-prime-2-September-23-2026}. \href{https://github.com/openai/math/blob/adc7f1241b42e322a6451854ab7e4b4c146bf78a/preprints/Fontaine-Mazur-modularity-at-the-prime-2-September-23-2026/paper.pdf}{PDF (GitHub)}; \href{https://github.com/openai/math/raw/adc7f1241b42e322a6451854ab7e4b4c146bf78a/preprints/Fontaine-Mazur-modularity-at-the-prime-2-September-23-2026/paper.pdf}{PDF (direct)}; \href{https://github.com/openai/math/blob/adc7f1241b42e322a6451854ab7e4b4c146bf78a/preprints/Fontaine-Mazur-modularity-at-the-prime-2-September-23-2026/README.md}{Manuscript page and citation}; \href{https://github.com/openai/math/blob/main/preprints/Fontaine-Mazur-modularity-at-the-prime-2-September-23-2026/paper.pdf}{Latest version}; Statement source: \href{https://github.com/openai/math/blob/adc7f1241b42e322a6451854ab7e4b4c146bf78a/preprints/Fontaine-Mazur-modularity-at-the-prime-2-September-23-2026/build/sections/introduction.tex\#L20}{thm:main}.
\item \href{https://github.com/openai/math/blob/adc7f1241b42e322a6451854ab7e4b4c146bf78a/CONTENTS.md}{OpenAI catalogue entry}.
\end{itemize}
\end{document}
